%
% dial_settlement.tex
%
% A Z specification of how a dial is ENDED and SETTLED: the reader slot, the
% listener's start and stop, the inline reader's last check and release, the
% no-reader settle, the Hub's own end and its bound.  It refines the two
% abstractions that docs/crash_attribution.tex makes (a "reader slot" that
% someone holds, a "settle" that someone performs) to the granularity of the
% code, because the final review rounds on the lux-ghhf branch found, again and
% again, a path on which an ended dial is never settled or is settled undrained
% (R7-F1, R8 R1, and the three defects of round 9).  This is the recurrence rule:
% the same class of defect in two or more rounds is a state-space problem, and
% the state space is finite.
%
% FIDELITY CONTROL writer_owe_refused: dial_settlement.tex with ONE change, in the
% section "The writer": WOweOk is refused after a Hub request or a recorded end
% (owe_ack consults `closing' through _require_live), by the added WOweRefused, so a
% scene written to the wire is not in the ledger of a current dial.  WRec is the
% design's.  Goal (S6):
%   wire = fset & ledger = fclear & lk = fclear & current = fset
% The goal s5bad = fset is NOT reached by this control alone (the refusal under a
% Hub request records nothing, and under a recorded end finds the cause set): S5
% needs this change together with a classification that ignores a Hub request,
% which is the control writer_records_display.  The rest of this header is that of
% the design file.
%
% THIS FILE IS THE PROPOSED DESIGN (round 4: round 3 plus the WRITER).  Round 3
% abstracted the writer away; the round-10 review found that the writer path breaks
% S5 and A2, and the section "The writer" models it.  Each fidelity control differs
% from it in the stated place only, and reproduces one trace.  The controls listed
% first below were committed against the round-3 file, which had no writer; they are
% one-point mutations of THAT file and are not regenerated here.  They still
% stand: with the writer idle (pc(aWr) = pidle, lk = fclear) every round-3 trace is
% a round-4 trace, the actor sets only grew, and the flags never clear, so each
% control's violation remains reachable in this file.  The writer controls are
% mutations of this one:
%   dial_settlement_writer_owe_refused_buggy.tex      the owed ack is refused after a
%                                                     Hub request or a recorded end (S6)
%   dial_settlement_writer_records_display_buggy.tex  today's code, TWO changes: the
%                                                     refused ack and a classification
%                                                     that ignores a Hub request (S5)
%   dial_settlement_drain_gives_up_buggy.tex          a drain that finds the socket
%                                                     lock held reports nothing more
%   dial_settlement_drain_no_late_buggy.tex           no timed-out branch: a settler
%                                                     with no asker waits for a stuck
%                                                     writer for ever (deadlock)
% In docs/:
%   dial_settlement_start_fail_buggy.tex       start_on on an ended dial keeps the slot
%   dial_settlement_stop_skip_drain_buggy.tex  the listener's exit reads ended twice
%                                              (a substitution of its steps, not a
%                                              one-point mutation)
%   dial_settlement_inline_strand_buggy.tex    a refused settle leaves no debt
% Also in docs/:
%   dial_settlement_split_settle_buggy.tex     the post-drain record is not first-wins
%   dial_settlement_atomic_refuse_buggy.tex    refusal and debt are two steps
%   dial_settlement_atomic_release_buggy.tex   release and debt-read are two steps
%   dial_settlement_defer_only_buggy.tex       a Hub request with no bound (T15 absent)
%   dial_settlement_timer_hub_only_buggy.tex   only a Hub request arms the bound
%   dial_settlement_s3_unguarded_buggy.tex     a capture without the current guard
%   dial_settlement_s4_overwrite_buggy.tex     a recorder that overwrites
%   dial_settlement_hub_record_buggy.tex       the Hub end records hub at once (S5)
%   dial_settlement_hub_peek_buggy.tex         the peek rule, an ack ahead of the FIN
% The code it governs: src/punt_lux/domain/hub/{display_connection,dial_reader,
% display_listener,listener_thread,display_replies,reply_sink,dial_generation,
% reader_slot,dial_drain,settler,bounded_ask,dial_end,dial_socket}.py.  Companions: crash_attribution.tex (what a capture
% charges), listener_locks.tex (the lock order rq < sl < st < lk < qu).
%
% Type-check:  fuzz docs/dial_settlement.tex
% Model-check: see the section Verification (one run at a time; each is a few
%              seconds and well under a gigabyte).
%
\documentclass[a4paper,10pt,fleqn]{article}
\usepackage[margin=1in]{geometry}
\usepackage{fuzz}

\begin{document}

\title{Dial Settlement: a Z Specification}
\author{Who ends a dial, who settles it, and who may let go}
\date{October 2026}
\maketitle

\section{What is modelled}

A \emph{dial} is one generation of the Hub's connection to the display.  It is
\emph{published} (current) from its birth until some one party \emph{captures}
it: takes its ledger, queues its death or carries its suspects, and makes it
not current.  Between the first end and the capture the dial is ended and
unsettled.  Five properties are claimed of that interval.

\begin{description}
  \item[S1] No dial stays ended (or end-requested), current and unsettled for
    ever.  Two forms.  \emph{S1a}, with a holder that makes progress: an end is
    always followed by a settle, by a party that was already running, with no
    timer.  \emph{S1b}, with a holder that never makes progress: the bound of
    rule~6 settles it.  A dial that is wedged is one that redial cannot replace,
    because it is still current.
  \item[S2] A dial ended by a \emph{dead} display is captured only after a drain
    that began after the end was recorded.  (The model's \texttt{cdisp} is
    always a dead display; see below for the live display and the carve-outs.)
  \item[S3] Exactly one capture per dial.  \emph{Structural:} every capture is
    guarded by \texttt{current} and clears it, so the flag cannot be set; the
    content of S3 is the guard, and the control \texttt{s3\_unguarded} removes it
    and trips the flag.
  \item[S4] The cause that is captured is the first one recorded.
    \emph{Structural:} \texttt{ended} has one first-wins writer; the control
    \texttt{s4\_overwrite} lets a recorder overwrite and trips the flag.
  \item[S5] No recorded cause contradicts the observation.  A display that was
    already dead when the settler observed the socket is never recorded as a Hub
    end; and (round 4) a writer never records a display end for a display that is
    alive, which is what a write that completed after a Hub request, and was
    then refused its ledger entry, did.  \emph{Structural in the design} (the
    cause is a function of the observation); the controls \texttt{hub\_record},
    \texttt{hub\_peek} and \texttt{writer\_records\_display} trip it.  The
    control needs two defects together, the refused ack and the classification
    that ignores a pending Hub request: the refused ack alone
    (\texttt{writer\_owe\_refused}) does not reach it, and the classification
    alone reaches nothing, because the design's writer records a display end
    only for a failure that the display caused.  The classification (rule~7b) is
    therefore \emph{policy}, stated to match the order \texttt{decide} uses (a
    Hub request outranks what a writer reports), and no checked property
    depends on it.
  \item[S6] (round 4) A scene whose bytes are on the wire of a current dial is in
    the ledger once its writer has let go of the socket lock; an end never
    withdraws that entry, and only a capture, which takes the ledger, may precede
    it.  The scene on the wire is owed.  Not a flag: a state predicate, checked as
    unreachability of \texttt{wire = fset $\land$ ledger = fclear $\land$
    lk = fclear $\land$ current = fset}.  The controls
    \texttt{writer\_owe\_refused} and \texttt{writer\_records\_display} reach it,
    each on its own.  A capture that lands before the entry is the exception to
    A1 below (\texttt{a1late}).
\end{description}

S2 is read with one more clause (round 4): a drain that finds the socket
readable never reports ``nothing more''.  This is \emph{structural}: its content
is \texttt{SDrain}'s guard, which is not enabled behind a held lock, and the flag
\texttt{gu} is identically clear in the design.  The control
\texttt{drain\_gives\_up} sets \texttt{gu} and trips \texttt{s2bad}.  What
replaces the give-up in the design is a timed-out branch, \texttt{SDrainLate},
which is checked: it ends in the bound's undrained capture, flagged
\texttt{s2lat}.

S1 is checked as reachability of a wedge state and by the deadlock check; S2 to
S5 are flags set by the operation that would violate them, and the goal is a
state with a flag set.  The reason is the one \texttt{listener\_locks.tex}
gives: ProB treats the primed state schema as part of each operation's guard, so
a violation written as a conjunct disables the operation instead of reporting
it.

\paragraph{The carve-outs of S2, with their bounds.}
Three paths capture a display or unresponsive end without the drain, because the
drain is not possible or not safe.  Each sets its own flag, never \texttt{s2bad}.
\begin{enumerate}
  \item \emph{Decode errors} (\texttt{s2dec}).  A frame that broke the stream
    (the reader's malformed-frame exit, \texttt{end\_stream(display)} in an
    inline read) cannot be read past.  The acknowledgements behind the bad frame
    are unread, and their scenes may be charged.  Bound: the acks in flight on the
    one dial, once.
  \item \emph{The bound of rule~6} (\texttt{s2blk}).  The Hub, tired of waiting
    for a holder that has not settled, captures without draining, whether the
    recorded cause is \texttt{display} or \texttt{unresponsive}.  Bound: every
    acknowledgement queued while the holder stalls, including what sits
    undelivered in the holder's frame buffer, and only after the bound (the stop
    join, two seconds).  A bound capture of a display end may therefore charge
    scenes whose acknowledgements were never read.  That is an accepted
    over-charge (DES-102).
  \item \emph{A live display} (unresponsive ends, for the holder's settle).  A
    display that is alive may send an acknowledgement after any drain, so no
    drain is ever complete.  S2 for an unresponsive end, settled by a holder, is
    that a drain that began after the record precedes the capture; an ack that
    arrives after that drain is bounded by one drain-to-capture gap.  This holds
    for the holder's settle only; the bound's capture of an unresponsive end is
    carve-out~2.  A dead display has its end of file queued behind its last byte
    on a same-host socket, so its drain is complete there; the claim that a drain
    is complete is restricted to same-host sockets, and that restriction is the
    contract (see \emph{What is not modelled}).  The classification of a
    send timeout as the display's end (a live display charged as dead) is a
    separate bound, B3, and is not modelled here.
\end{enumerate}

\section{The design}

Settlement belongs to the \emph{holder of the reader slot}; everyone else only
\emph{asks}.

\begin{enumerate}
  \item \textbf{Capture authority is the slot.}  Only a holder of the slot
    (the listener thread, an inline reader, or a releaser taking over) captures a
    dial.  In code: \texttt{close\_if\_current} leaves the connection's
    interface; the slot hands its holder a \texttt{Settler}, the one object with
    a \texttt{settle()} method, and the listener, the inline reader and the
    no-reader settle are the only holders.  There is no flag with a default
    (\texttt{by\_reader=True}) and no Hub door.  The model's one operation that
    captures without holding the slot is the bound of rule~6, and it is
    labelled.
  \item \textbf{Ends record or request; they never capture.}  A display or
    unresponsive end records its cause, first wins.  A Hub end sets a request
    flag (\texttt{hubreq}) and records no cause.  Both then \emph{ask for a
    settle}: take the slot if it is free and settle at once; if it is held, leave
    a debt on the slot and wait for the bound.
  \item \textbf{Whoever lets go of the slot honours the debt.}  A release
    (listener start failure, inline exit) that finds a debt keeps the slot,
    settles, and looks again.  Every path that gives up the reader role releases
    the slot this way, including a start that raises because the dial was ended.
  \item \textbf{Settle is two steps, a drain and a decision.}  The drain reads
    the socket to the point where it would block and notes whether end of file
    came.  The decision fixes the cause with a first-wins record: the observed
    end of file is \texttt{display}, else a pending Hub request is \texttt{hub};
    a cause recorded earlier stands.  If the record stands and was made after the
    drain began, the drain was not complete: \emph{the settler drains again}
    before it captures.  After the record is fixed nothing can change the cause,
    so the second drain is the last.  Both settlers, the listener and the
    releaser, run these steps.
  \item \textbf{A Hub request acts at once on what it can.}  It refuses writes
    and fails waiters with the Hub's error immediately, without fixing a cause.
    It does not capture, and it does not redial.
  \item \textbf{The wait is bounded (T15).}  An asker refused behind a held slot
    waits.  If the holder has not settled within the bound (the stop join,
    2~s), the asker records its cause (a Hub request records \texttt{hub}) with
    the same first-wins record, and captures.  A capture of cause \texttt{hub}
    charges nothing; its scenes are carried.  The holder's later settle finds the
    dial not current and does nothing.  The bound applies to \emph{every} asker,
    not only to a Hub request: a display end discovered behind a holder that
    never looks is otherwise wedged for as long as the holder stalls
    (the control \texttt{timer\_hub\_only}).
  \item \textbf{The writer (round 4).}  A write holds the socket lock across
    \texttt{sendall} and the owed-ack entry.  Three consequences.  (a)~Once the
    bytes are on the wire, the owed ack is recorded whatever the end state: a Hub
    request or a recorded cause does not refuse it; only a dial already captured
    does.  (b)~A write that fails after a Hub request has been made is the
    Hub's failure, not the display's: the writer records no cause and leaves the
    decision to the settler, whose drain finds a display that really died.
    (c)~A drain that finds the socket readable and the lock held by a writer
    \emph{waits} for the writer; it does not report ``nothing more''.  The wait
    has a ceiling, twice the one-second send timeout
    (\texttt{SEND\_TIMEOUT\_CEILING\_SECONDS}, two seconds).  A writer that fails or times out
    releases the lock within it.  A writer that keeps making partial progress
    (\texttt{WStall}) does not, and rule~6 does \emph{not} bound that case,
    because it is armed by a refused asker and a settler that took the slot
    itself has none.  At the ceiling the drain reports itself incomplete
    (\texttt{DrainIncompleteError}) and the settler hands its capture to the
    bound, \texttt{ReaderSlot.\_late}: it waits the bound, two seconds, for
    anyone else to settle the dial, then records its cause and captures
    undrained (\texttt{SDrainLate}, \texttt{LateCapture}).  The settle therefore
    ends within the ceiling plus the bound, four seconds, with no asker
    required; the capture shuts the socket, which is what
    wakes the stuck writer (\texttt{WStuckOut}).  The obligation this places on
    the code: a capture shuts the socket whatever path captured, so a TLS
    \texttt{sendall} wedged on a display that reads nothing returns within that
    bound (the test \texttt{test\_retiring\_a\_dial\_releases\_a\_sender\_wedged\_on\_a\_tls\_socket}).
    The late capture is the fifth exception to A2 (\texttt{s2lat}), with
    the bound of rule~6 plus the ceiling.  The lock order
    $rq < sl < st < lk < qu$ is kept: the waiting settler holds no lock that the
    writer takes, and the writer, which takes \texttt{qu} for the ack under
    \texttt{lk}, waits on nothing the settler holds.  The alternative that was
    weighed in round 1, a drain that returns ``incomplete'' so that the settler
    drains again, is the ceiling taken in polling steps; the code does not poll,
    it hands the capture to the bound.
\end{enumerate}

Rules 4 and 5 are what close S5.  The model assumes the observation is one
linearization point (the drain to \texttt{EAGAIN} together with the end-of-file
test).  A display that dies after it dies after the Hub's end; that is not a
defect, and no rule can distinguish it.  The residual widens to the bound if
rule~6 fires: the Hub records \texttt{hub} for a display it has not observed, and
the flag \texttt{s5blk} says so.

\paragraph{The three atomicity obligations.}
S1a holds only if each of these is one critical section of the slot's condition
lock; each has a control that splits it and fails.  With the timer present a
broken obligation does not wedge: it degrades to the bound of rule~6, that is to
the two-second wait and an undrained capture (\texttt{s2blk}).
\begin{enumerate}
  \item \textbf{Refusal and debt}: observing the slot held and setting the debt
    (\texttt{atomic\_refuse}).
  \item \textbf{Release and debt-read}: reading the debt and freeing the slot
    (\texttt{atomic\_release}).
  \item \textbf{The post-drain record}: the first-wins record of the decided
    cause, with the re-drain when it is lost (\texttt{split\_settle}).  This one
    is the dial's \texttt{end()} under the sink's own lock, not the slot's.
\end{enumerate}

\section{Carriers and bounds}

One dial.  Every operation in the code is scoped to a generation: a reader, a
slot, a write failure, and a capture all name the dial they act on, and the Hub
redials only when the dial is no longer current.  A second dial therefore runs
in a disjoint product with the first, and cannot interact with it; the model
includes \texttt{Redial}, enabled only on an unpublished dial, to say so.  Six
actors, each of which acts at most once: the starter (\texttt{start\_on}), the
stopper (a Hub end: \texttt{stop} or \texttt{close}), a discoverer of a display
or unresponsive end (the keepalive, or a write failure seen as such), one inline
reader, the listener thread, and (round 4) the writer of one scene.  Socket content is zero or one unread acknowledgement,
which is enough to tell a drained capture from an undrained one; it arrives at
any time while the display lives.  The dial is born holding the listener's slot
(the registry's \texttt{open\_listening}) or free (a library link).  A holder may
\emph{block} (a handler on the listener thread, an inline read that never
returns): it then takes none of its own steps.

A send timeout that does not abort a wedged \texttt{sendall} is modelled as the
stuck writer (\texttt{WStall}): the writer keeps the socket lock until a capture
shuts the socket.  The ceiling on the drain's lock wait and the bound that
follows it (rule~7c) are what end such a settle.

\begin{zed}
  CAUSE ::= nocause | cdisp | cunresp | chub
\end{zed}

\begin{zed}
  SLOT ::= sfree | slistener | sinline
\end{zed}

\begin{zed}
  FLAG ::= fclear | fset
\end{zed}

\begin{zed}
  ACTOR ::= aSt | aSp | aDi | aIn | aRd | aWr
\end{zed}

The program counters.  \texttt{pst0}: the starter has not yet started the
listener.  \texttt{prel0}: it failed and owes a release.  \texttt{phub}: a Hub
end is to be requested.  \texttt{pend}: a display end is to be recorded.
\texttt{psu}: the no-reader settle is to be tried.  \texttt{pwait}: refused, and
waiting for the bound.  \texttt{pset}: this actor holds the slot and owes a
drain.  \texttt{pdec}: it has drained and owes the decision.  \texttt{prel}: it
is about to release.  \texttt{iread}, \texttt{ilast}: the inline reader is
reading, and has finished reading and will make its last check.  \texttt{rtop}:
the listener is at the top of its loop.  \texttt{rnone}: no listener exists.
The writer's: \texttt{wsend}, it holds the socket lock with the scene not yet on
the wire; \texttt{wowe}, the bytes are on the wire and the ack is to be owed;
\texttt{wunl}, the entry is made and the lock is to be let go; \texttt{wfail},
the lock is released after a failure and the failure is to be classified.

\begin{zed}
  PC ::= pidle | pst0 | prel0 | phub | pend | psu | pwait | pset | pdec | prel |
    pdone | iread | ilast | rtop | rnone | rdone | wsend | wowe | wunl | wfail |
    wstuck | plate
\end{zed}

\texttt{wstuck}: the writer holds the socket lock and never lets go of it on its
own (a TLS \texttt{sendall} that makes slow partial progress, so that its send
timeout never fires).  \texttt{plate}: a settler whose drain could not be made
(the lock was held past the ceiling) has handed its capture to the bound.

\begin{schema}{Dial}
  ended : CAUSE \\
  slot : SLOT \\
  debt : FLAG \\
  current : FLAG \\
  sock : 0 \upto 1 \\
  capN : 0 \upto 2 \\
  hubreq : FLAG
\end{schema}

\begin{schema}{Env}
  launched : FLAG \\
  eof : FLAG \\
  redialed : FLAG \\
  blocked : FLAG
\end{schema}

\begin{schema}{Flags}
  s2bad : FLAG \\
  s3bad : FLAG \\
  s4bad : FLAG \\
  s5bad : FLAG
\end{schema}

\begin{schema}{Carve}
  s2blk : FLAG \\
  s2dec : FLAG \\
  s5blk : FLAG \\
  s2lat : FLAG \\
  a1late : FLAG
\end{schema}

\begin{schema}{Pc}
  pc : ACTOR \fun PC
\end{schema}

\begin{schema}{Dc}
  dcs : CAUSE
\end{schema}

\begin{schema}{RLoc}
  rcause : CAUSE \\
  rdec : FLAG \\
  rdb : FLAG
\end{schema}

\begin{schema}{Loc}
  obs : ACTOR \fun CAUSE \\
  dsaw : ACTOR \fun CAUSE \\
  gu : FLAG
\end{schema}

\begin{schema}{Wr}
  lk : FLAG \\
  wire : FLAG \\
  ledger : FLAG \\
  wcs : CAUSE
\end{schema}

\texttt{lk} is the socket lock held by the writer (the reader's \texttt{recv}
also takes it, but the reader is a settler's own business and is not modelled as
a lock user).  \texttt{wire}: the scene's bytes are on the wire.  \texttt{ledger}:
the scene is owed an ack in the dial's ledger.  \texttt{wcs}: the cause the
writer's failure reports.  \texttt{gu}: the last drain gave up with the lock held;
the design never sets it, and the control \texttt{drain\_gives\_up} does.

\texttt{ended} is the first recorded cause.  \texttt{hubreq} is a request, not a
cause.  \texttt{eof} says that the display has died and its end of file is
queued behind whatever is unread; \texttt{sock} is that unread
acknowledgement.  \texttt{sock} stands for the socket \emph{and} the reader's
undelivered buffer: that is sound only because rule~1 makes the settler the sole
owner of the frame buffer, so what the reader has taken and not delivered is the
settler's own, and a drain by the settler leaves neither (the gluing is in the
section Refinement).  \texttt{dcs} is the discoverer's cause.  \texttt{obs} is
what an actor's last drain saw (end of file, or nothing); \texttt{dsaw} is the
cause that stood when that drain began.  \texttt{rcause}, \texttt{rdec} and
\texttt{rdb} are locals that only controls use.

\begin{schema}{Sys}
  Dial \\ Env \\ Flags \\ Carve \\ Pc \\ Dc \\ RLoc \\ Loc \\ Wr
\end{schema}

\begin{axdef}
  nxt : ACTOR \fun PC
\where
  nxt = \{ aSt \mapsto prel, aSp \mapsto prel, aDi \mapsto prel, aIn \mapsto prel,
    aRd \mapsto rdone, aWr \mapsto prel \}
\end{axdef}

Where an actor goes after its settle: a caller to its release, the listener to
its end.

\begin{schema}{Init}
  Sys'
\where
  ended' = nocause \\
  debt' = fclear \\
  current' = fset \\
  sock' \in 0 \upto 1 \\
  capN' = 0 \\
  hubreq' = fclear \\
  launched' = fclear \\
  eof' = fclear \\
  redialed' = fclear \\
  blocked' = fclear \\
  s2bad' = fclear \\
  s3bad' = fclear \\
  s4bad' = fclear \\
  s5bad' = fclear \\
  s2blk' = fclear \\
  s2dec' = fclear \\
  s5blk' = fclear \\
  s2lat' = fclear \\
  a1late' = fclear \\
  dcs' = cdisp \\
  rcause' = nocause \\
  rdec' = fclear \\
  rdb' = fclear \\
  obs' = ACTOR \cross \{ nocause \} \\
  dsaw' = ACTOR \cross \{ nocause \} \\
  gu' = fclear \\
  lk' = fclear \\
  wire' = fclear \\
  ledger' = fclear \\
  wcs' = nocause \\
  (slot' = slistener \land pc' = \{ aSt \mapsto pst0, aSp \mapsto pidle,
    aDi \mapsto pidle, aIn \mapsto pidle, aRd \mapsto rnone, aWr \mapsto pidle \}) \lor
  (slot' = sfree \land pc' = \{ aSt \mapsto pdone, aSp \mapsto pidle,
    aDi \mapsto pidle, aIn \mapsto pidle, aRd \mapsto rnone, aWr \mapsto pidle \})
\end{schema}

The frames.  Each operation names, through one of these, which parts of the
state it may change; the rest are unchanged.

\begin{schema}{OnlyPc}
  \Xi Dial \\ \Xi Env \\ \Xi Flags \\ \Xi Carve \\ \Delta Pc \\ \Xi Dc \\
  \Xi RLoc \\ \Xi Loc \\ \Xi Wr
\end{schema}

\begin{schema}{DialPc}
  \Delta Dial \\ \Xi Env \\ \Xi Flags \\ \Xi Carve \\ \Delta Pc \\ \Xi Dc \\
  \Xi RLoc \\ \Xi Loc \\ \Xi Wr
\end{schema}

\begin{schema}{DialOnly}
  \Delta Dial \\ \Xi Env \\ \Xi Flags \\ \Xi Carve \\ \Xi Pc \\ \Xi Dc \\
  \Xi RLoc \\ \Xi Loc \\ \Xi Wr
\end{schema}

\begin{schema}{EnvOnly}
  \Xi Dial \\ \Delta Env \\ \Xi Flags \\ \Xi Carve \\ \Xi Pc \\ \Xi Dc \\
  \Xi RLoc \\ \Xi Loc \\ \Xi Wr
\end{schema}

\begin{schema}{EnvPc}
  \Xi Dial \\ \Delta Env \\ \Xi Flags \\ \Xi Carve \\ \Delta Pc \\ \Xi Dc \\
  \Xi RLoc \\ \Xi Loc \\ \Xi Wr
\end{schema}

\begin{schema}{PcDc}
  \Xi Dial \\ \Xi Env \\ \Xi Flags \\ \Xi Carve \\ \Delta Pc \\ \Delta Dc \\
  \Xi RLoc \\ \Xi Loc \\ \Xi Wr
\end{schema}

\begin{schema}{Drain}
  \Delta Dial \\ \Xi Env \\ \Xi Flags \\ \Xi Carve \\ \Delta Pc \\ \Xi Dc \\
  \Xi RLoc \\ \Delta Loc \\ \Xi Wr
\end{schema}

\begin{schema}{Cap}
  \Delta Dial \\ \Xi Env \\ \Delta Flags \\ \Xi Carve \\ \Delta Pc \\ \Xi Dc \\
  \Xi RLoc \\ \Xi Loc \\ \Xi Wr
\end{schema}

\begin{schema}{CapCarve}
  \Delta Dial \\ \Xi Env \\ \Delta Flags \\ \Delta Carve \\ \Delta Pc \\ \Xi Dc \\
  \Xi RLoc \\ \Xi Loc \\ \Xi Wr
\end{schema}

\begin{schema}{WrPc}
  \Xi Dial \\ \Xi Env \\ \Xi Flags \\ \Xi Carve \\ \Delta Pc \\ \Xi Dc \\
  \Xi RLoc \\ \Xi Loc \\ \Delta Wr
\end{schema}

\begin{schema}{WrLate}
  \Xi Dial \\ \Xi Env \\ \Xi Flags \\ \Delta Carve \\ \Delta Pc \\ \Xi Dc \\
  \Xi RLoc \\ \Xi Loc \\ \Delta Wr
\end{schema}

\begin{schema}{WrRec}
  \Delta Dial \\ \Xi Env \\ \Delta Flags \\ \Xi Carve \\ \Delta Pc \\ \Xi Dc \\
  \Xi RLoc \\ \Xi Loc \\ \Xi Wr
\end{schema}

\section{The environment}

\begin{schema}{Die}
  EnvOnly
\where
  eof = fclear \\
  eof' = fset \\
  launched' = launched \\
  redialed' = redialed \\
  blocked' = blocked
\end{schema}

The display dies.  Its end of file is queued behind any acknowledgement it has
already sent.

\begin{schema}{Ack}
  DialOnly
\where
  current = fset \\
  eof = fclear \\
  sock = 0 \\
  sock' = 1 \\
  ended' = ended \\
  slot' = slot \\
  debt' = debt \\
  current' = current \\
  capN' = capN \\
  hubreq' = hubreq
\end{schema}

A living display sends an acknowledgement, at any time.  This is what the round-1
model lacked: without it no capture could ever race an ack.

\begin{schema}{Block}
  EnvOnly
\where
  blocked = fclear \\
  blocked' = fset \\
  launched' = launched \\
  eof' = eof \\
  redialed' = redialed
\end{schema}

The holder of the slot stops making progress, for good: a handler that never
returns.  It takes none of the steps that need it (\texttt{RNotice},
\texttt{IReadIdle}, \texttt{IReadNotice}, \texttt{IBad}, \texttt{RBad}).  Steps
that a holder takes once it has begun to settle are not blocked: the drain does
not wait on the display.

\begin{schema}{Redial}
  EnvOnly
\where
  current = fclear \\
  redialed = fclear \\
  redialed' = fset \\
  launched' = launched \\
  eof' = eof \\
  blocked' = blocked
\end{schema}

The Hub dials again, and only on a dial that is no longer current.  A dial
that is ended but still current is never replaced: that is the wedge.

\begin{schema}{Halt}
  \Xi Sys
\where
  current = fclear \lor (ended = nocause \land hubreq = fclear)
\end{schema}

The model may stop only when the dial is settled or has not been ended.  A
state with the dial ended or end-requested, still current, and nothing enabled
is a deadlock, and probcli reports it.

\section{The starter}

\begin{schema}{StOk}
  EnvPc
\where
  pc~aSt = pst0 \\
  current = fset \\
  ended = nocause \\
  hubreq = fclear \\
  launched' = fset \\
  eof' = eof \\
  redialed' = redialed \\
  blocked' = blocked \\
  pc' = pc \oplus \{ aSt \mapsto pdone, aRd \mapsto rtop \}
\end{schema}

\texttt{start\_on}: the dial is not closed, and the thread starts.

\begin{schema}{StFail}
  OnlyPc
\where
  pc~aSt = pst0 \\
  pc' = pc \oplus \{ aSt \mapsto prel0 \}
\end{schema}

\texttt{start\_on} fails: the dial was found ended or requested
(\texttt{ConnectionClosedError}), or the thread could not be started.  The
design keeps this unguarded, which is the stronger claim: it releases on both.
The starter owes a release of the slot it was born holding.

\begin{schema}{StRel0}
  DialPc
\where
  pc~aSt = prel0 \\
  slot = slistener \\
  ended' = ended \\
  current' = current \\
  sock' = sock \\
  capN' = capN \\
  hubreq' = hubreq \\
  (debt = fset \land debt' = fclear \land slot' = sinline \land
    pc' = pc \oplus \{ aSt \mapsto pset \}) \lor
  (debt = fclear \land debt' = debt \land slot' = sfree \land
    pc' = pc \oplus \{ aSt \mapsto phub \})
\end{schema}

The release.  With a debt the starter keeps the slot and settles; without, it
frees the slot and goes on to request the Hub's end of the dial it could not
start (\texttt{retire(born, hub)}).

\section{The ends}

\begin{schema}{DStartDisp}
  PcDc
\where
  pc~aDi = pidle \\
  dcs' = cdisp \\
  pc' = pc \oplus \{ aDi \mapsto pend \}
\end{schema}

\begin{schema}{DStartUnresp}
  PcDc
\where
  pc~aDi = pidle \\
  dcs' = cunresp \\
  pc' = pc \oplus \{ aDi \mapsto pend \}
\end{schema}

A discoverer: a write that failed with a broken pipe (the display's end, so the
display is dead), or the keepalive giving up (unresponsive, the display perhaps
alive).

\begin{schema}{RecEnd}
  DialPc
\where
  pc~aDi = pend \\
  current = fset \\
  (dcs = cunresp \lor eof = fset) \\
  slot' = slot \\
  debt' = debt \\
  current' = current \\
  sock' = sock \\
  capN' = capN \\
  hubreq' = hubreq \\
  ((ended = nocause \land ended' = dcs) \lor (ended \neq nocause \land ended' = ended)) \\
  pc' = pc \oplus \{ aDi \mapsto psu \}
\end{schema}

\begin{schema}{RecGone}
  OnlyPc
\where
  pc~aDi = pend \\
  current = fclear \\
  pc' = pc \oplus \{ aDi \mapsto pdone \}
\end{schema}

The cause is recorded, first wins.  A dial no longer current is left alone.

\begin{schema}{SPStart}
  OnlyPc
\where
  pc~aSp = pidle \\
  pc' = pc \oplus \{ aSp \mapsto phub \}
\end{schema}

The stopper: \texttt{stop} (when a listener exists) or \texttt{close}.  Both end
by asking for the Hub's end.  A \texttt{stop} on a dial with no listener
returns without ending anything; that is a subset of \texttt{close}.

\begin{schema}{HubReq}
  DialPc \\
  a? : ACTOR
\where
  a? \in \{ aSt, aSp \} \\
  pc~a? = phub \\
  current = fset \\
  ended' = ended \\
  slot' = slot \\
  debt' = debt \\
  current' = current \\
  sock' = sock \\
  capN' = capN \\
  hubreq' = fset \\
  pc' = pc \oplus \{ a? \mapsto psu \}
\end{schema}

\begin{schema}{HubGone}
  OnlyPc \\
  a? : ACTOR
\where
  a? \in \{ aSt, aSp \} \\
  pc~a? = phub \\
  current = fclear \\
  pc' = pc \oplus \{ a? \mapsto pdone \}
\end{schema}

The Hub's end sets a request.  The request itself is what refuses writes and
fails waiters at once (rule~5); it records no cause and captures nothing.  That
refusal of writes and failing of waiters is a stated rule, not modelled state:
no checked property depends on it.

\section{The no-reader settle, the wait and the release}

\begin{schema}{SUGone}
  OnlyPc \\
  a? : ACTOR
\where
  a? \in \{ aSt, aSp, aDi, aWr \} \\
  pc~a? = psu \\
  current = fclear \\
  pc' = pc \oplus \{ a? \mapsto pdone \}
\end{schema}

\begin{schema}{SUGrant}
  DialPc \\
  a? : ACTOR
\where
  a? \in \{ aSt, aSp, aDi, aWr \} \\
  pc~a? = psu \\
  current = fset \\
  slot = sfree \\
  slot' = sinline \\
  ended' = ended \\
  debt' = debt \\
  current' = current \\
  sock' = sock \\
  capN' = capN \\
  hubreq' = hubreq \\
  pc' = pc \oplus \{ a? \mapsto pset \}
\end{schema}

\begin{schema}{SURefuse}
  DialPc \\
  a? : ACTOR
\where
  a? \in \{ aSt, aSp, aDi, aWr \} \\
  pc~a? = psu \\
  current = fset \\
  slot \neq sfree \\
  debt' = fset \\
  ended' = ended \\
  slot' = slot \\
  current' = current \\
  sock' = sock \\
  capN' = capN \\
  hubreq' = hubreq \\
  pc' = pc \oplus \{ a? \mapsto pwait \}
\end{schema}

The slot's ask (\texttt{ReaderSlot.ask\_settle}).  A free slot is taken and the
caller settles through its \texttt{Settler}.  A held
slot is left a debt, in the same step as the refusal: the holder cannot release
between the refusal and the debt.  The asker then waits.

\begin{schema}{WaitGone}
  OnlyPc \\
  a? : ACTOR
\where
  a? \in \{ aSt, aSp, aDi, aWr \} \\
  pc~a? = pwait \\
  current = fclear \\
  pc' = pc \oplus \{ a? \mapsto pdone \}
\end{schema}

\begin{schema}{Timeout}
  CapCarve \\
  a? : ACTOR
\where
  a? \in \{ aSt, aSp, aDi, aWr \} \\
  pc~a? = pwait \\
  current = fset \\
  ((ended = nocause \land ended' = chub) \lor (ended \neq nocause \land ended' = ended)) \\
  current' = fclear \\
  capN' = capN + 1 \\
  slot' = slot \\
  debt' = debt \\
  sock' = sock \\
  hubreq' = hubreq \\
  s2bad' = s2bad \\
  ((capN \geq 1 \land s3bad' = fset) \lor (capN = 0 \land s3bad' = s3bad)) \\
  s4bad' = s4bad \\
  s5bad' = s5bad \\
  ((ended' \in \{ cdisp, cunresp \} \land sock > 0 \land s2blk' = fset) \lor
    (\lnot (ended' \in \{ cdisp, cunresp \} \land sock > 0) \land s2blk' = s2blk)) \\
  s2dec' = s2dec \\
  ((ended = nocause \land eof = fset \land s5blk' = fset) \lor
    (\lnot (ended = nocause \land eof = fset) \land s5blk' = s5blk)) \\
  s2lat' = s2lat \\
  a1late' = a1late \\
  pc' = pc \oplus \{ a? \mapsto pdone \}
\end{schema}

Rule~6: the bound has passed.  The asker records its cause with the first-wins
record and captures, though it does not hold the slot and does not drain.  It is
the one capture that is not by a holder.  A \texttt{hub} capture charges nothing.
The two flags say what this costs: a display end captured with an
acknowledgement unread (\texttt{s2blk}), and a \texttt{hub} recorded for a
display that has died unobserved (\texttt{s5blk}).

\begin{schema}{SDrain}
  Drain \\
  a? : ACTOR
\where
  pc~a? = pset \\
  slot \neq sfree \\
  current = fset \\
  (ended \neq nocause \lor hubreq = fset \lor eof = fset) \\
  (lk = fclear \lor (sock = 0 \land eof = fclear)) \\
  gu' = fclear \\
  sock' = 0 \\
  ended' = ended \\
  slot' = slot \\
  debt' = debt \\
  current' = current \\
  capN' = capN \\
  hubreq' = hubreq \\
  ((eof = fset \land obs' = obs \oplus \{ a? \mapsto cdisp \}) \lor
    (eof = fclear \land obs' = obs \oplus \{ a? \mapsto nocause \})) \\
  dsaw' = dsaw \oplus \{ a? \mapsto ended \} \\
  pc' = pc \oplus \{ a? \mapsto pdec \}
\end{schema}

The drain: the socket is read to the point of blocking.  It notes whether end of
file came, and which cause stood when it began.  It is not enabled while a writer
holds the socket lock when the socket is readable (an ack, or the end of file):
the drain takes the lock only to read, so a socket that \texttt{select} reports
not readable is looked at without it, and the guard is
\texttt{lk = fclear} $\lor$ (\texttt{sock = 0} $\land$ \texttt{eof = fclear}).
Behind a held lock the settler waits for the writer, which is the environment's
step, or gives up at the ceiling, which is \texttt{SDrainLate}.  A drain that
took the lock-held branch and reported ``nothing more'' with the socket unread
is the control \texttt{drain\_gives\_up}.

\begin{schema}{SDecide}
  Cap \\
  a? : ACTOR
\where
  pc~a? = pdec \\
  slot \neq sfree \\
  current = fset \\
  ((ended = nocause \land obs~a? \neq nocause \land ended' = obs~a?) \lor
    (ended = nocause \land obs~a? = nocause \land hubreq = fset \land ended' = chub) \lor
    (ended \neq nocause \land ended' = ended \land
      \lnot (ended \in \{ cdisp, cunresp \} \land dsaw~a? = nocause))) \\
  sock' = sock \\
  current' = fclear \\
  capN' = capN + 1 \\
  slot' = slot \\
  debt' = debt \\
  hubreq' = hubreq \\
  ((sock > 0 \land (ended' = cdisp \lor (gu = fset \land ended' = cunresp)) \land
      s2bad' = fset) \lor
    (\lnot (sock > 0 \land (ended' = cdisp \lor (gu = fset \land ended' = cunresp))) \land
      s2bad' = s2bad)) \\
  ((capN \geq 1 \land s3bad' = fset) \lor (capN = 0 \land s3bad' = s3bad)) \\
  s4bad' = s4bad \\
  s5bad' = s5bad \\
  pc' = pc \oplus \{ a? \mapsto nxt~a? \}
\end{schema}

The decision, by a holder: the first-wins record of the cause, then the
capture.  The cause is the end of file the drain saw, else a pending Hub
request; a cause already recorded stands, unless it is a display or
unresponsive end recorded since the drain began, in which case this step is not
enabled and \texttt{SRedrain} is.

\begin{schema}{SRedrain}
  OnlyPc \\
  a? : ACTOR
\where
  pc~a? = pdec \\
  slot \neq sfree \\
  current = fset \\
  ended \in \{ cdisp, cunresp \} \\
  dsaw~a? = nocause \\
  pc' = pc \oplus \{ a? \mapsto pset \}
\end{schema}

The record was lost to, or made after, the drain: drain again.  After this the
record stands and nothing can change the cause, so this loop runs at most once.

\begin{schema}{SettleGone}
  OnlyPc \\
  a? : ACTOR
\where
  pc~a? \in \{ pset, pdec \} \\
  current = fclear \\
  pc' = pc \oplus \{ a? \mapsto nxt~a? \}
\end{schema}

The dial was captured by someone else meanwhile (the bound of rule~6).  The
holder does nothing.

\begin{schema}{SDrainLate}
  OnlyPc \\
  a? : ACTOR
\where
  pc~a? = pset \\
  slot \neq sfree \\
  current = fset \\
  (ended \neq nocause \lor hubreq = fset \lor eof = fset) \\
  lk = fset \\
  (sock > 0 \lor eof = fset) \\
  pc' = pc \oplus \{ a? \mapsto plate \}
\end{schema}

The drain could not be made: the socket is readable, and a writer has held the
socket lock past the ceiling (twice the send timeout).  The drain reports itself
incomplete (\texttt{DrainIncompleteError}), and the settler, which holds the slot
and needs no refused asker, hands its capture to the bound.  The step is enabled
whenever the lock is held and the socket readable, which over-approximates the
code (the ceiling may not have passed); an enabled step is a possible step, and
a model that permits more is the safer one.

\begin{schema}{LateCapture}
  CapCarve \\
  a? : ACTOR
\where
  pc~a? = plate \\
  current = fset \\
  ((ended = nocause \land ended' = chub) \lor (ended \neq nocause \land ended' = ended)) \\
  current' = fclear \\
  capN' = capN + 1 \\
  slot' = slot \\
  debt' = debt \\
  sock' = sock \\
  hubreq' = hubreq \\
  s2bad' = s2bad \\
  ((capN \geq 1 \land s3bad' = fset) \lor (capN = 0 \land s3bad' = s3bad)) \\
  s4bad' = s4bad \\
  s5bad' = s5bad \\
  s2blk' = s2blk \\
  s2dec' = s2dec \\
  ((ended' \in \{ cdisp, cunresp \} \land sock > 0 \land s2lat' = fset) \lor
    (\lnot (ended' \in \{ cdisp, cunresp \} \land sock > 0) \land s2lat' = s2lat)) \\
  ((ended = nocause \land eof = fset \land s5blk' = fset) \lor
    (\lnot (ended = nocause \land eof = fset) \land s5blk' = s5blk)) \\
  a1late' = a1late \\
  pc' = pc \oplus \{ a? \mapsto nxt~a? \}
\end{schema}

\begin{schema}{LateGone}
  OnlyPc \\
  a? : ACTOR
\where
  pc~a? = plate \\
  current = fclear \\
  pc' = pc \oplus \{ a? \mapsto nxt~a? \}
\end{schema}

\texttt{BoundedAsk.settle\_by\_bound}, called by the settler on itself
(\texttt{ReaderSlot.\_late}).  It waits the bound for the dial to be settled by
someone else (\texttt{LateGone}); past it, it records its cause, first wins,
\texttt{hub} when none stands, and captures without a drain.  This is the
bound's capture, taken by the holder: the fifth exception to A2, with its own
flag \texttt{s2lat} (a display or unresponsive end captured with an ack unread,
because the drain could not be made), and, as for every bound capture,
\texttt{s5blk} when \texttt{hub} is recorded for a display that has died.  The
capture is what releases the writer (\texttt{WStuckOut}): it shuts the socket.

\begin{schema}{RelOwes}
  DialPc \\
  a? : ACTOR
\where
  a? \in \{ aSt, aSp, aDi, aIn, aWr \} \\
  pc~a? = prel \\
  slot = sinline \\
  debt = fset \\
  debt' = fclear \\
  ended' = ended \\
  slot' = slot \\
  current' = current \\
  sock' = sock \\
  capN' = capN \\
  hubreq' = hubreq \\
  pc' = pc \oplus \{ a? \mapsto pset \}
\end{schema}

\begin{schema}{RelFree}
  DialPc \\
  a? : ACTOR
\where
  a? \in \{ aSt, aSp, aDi, aIn, aWr \} \\
  pc~a? = prel \\
  slot = sinline \\
  debt = fclear \\
  slot' = sfree \\
  ended' = ended \\
  debt' = debt \\
  current' = current \\
  sock' = sock \\
  capN' = capN \\
  hubreq' = hubreq \\
  pc' = pc \oplus \{ a? \mapsto pdone \}
\end{schema}

The release.  A debt, left by a refusal that came before it, is honoured by
keeping the slot and settling again.  A debt left after the release finds the
slot free.

\section{The writer}

One writer sends one scene (\texttt{send\_scene}).  Its steps are those of
\texttt{DialSocket.\_write\_locked}: take the socket lock and ask whether the
dial is open; put the bytes on the wire; record the ack the scene is owed; let go
of the lock; and only then, if anything raised, tell the dial
(\texttt{\_write\_failed}), which classifies the failure and asks for a settle.
The classification and the owed-ack entry are the two places where round 3's
abstraction hid a defect.

\begin{schema}{WLock}
  WrPc
\where
  pc~aWr = pidle \\
  lk = fclear \\
  current = fset \\
  ended = nocause \\
  hubreq = fclear \\
  lk' = fset \\
  wire' = wire \\
  ledger' = ledger \\
  wcs' = wcs \\
  pc' = pc \oplus \{ aWr \mapsto wsend \}
\end{schema}

The lock is free and the dial is open (\texttt{\_require\_open}): neither a
recorded end nor a Hub request.  Otherwise nothing is written and nothing is
ended.

\begin{schema}{WSendDone}
  WrPc
\where
  pc~aWr = wsend \\
  wire' = fset \\
  lk' = lk \\
  ledger' = ledger \\
  wcs' = wcs \\
  pc' = pc \oplus \{ aWr \mapsto wowe \}
\end{schema}

\begin{schema}{WSendDead}
  WrPc
\where
  pc~aWr = wsend \\
  eof = fset \\
  lk' = fclear \\
  wire' = wire \\
  ledger' = ledger \\
  wcs' = cdisp \\
  pc' = pc \oplus \{ aWr \mapsto wfail \}
\end{schema}

\begin{schema}{WSendTimeout}
  WrPc
\where
  pc~aWr = wsend \\
  lk' = fclear \\
  wire' = wire \\
  ledger' = ledger \\
  wcs' = cunresp \\
  pc' = pc \oplus \{ aWr \mapsto wfail \}
\end{schema}

The send fails: the display is dead (a broken pipe, the end of file queued
behind it), or the send timeout fires on a display that is alive and reads
nothing.  The lock is released by the failure, before the dial is told.  The
timeout is what bounds a writer in \texttt{sendall}; it is always enabled, which
is the assumption that a writer making no progress for the send timeout fails.
A writer that makes slow partial progress is not bounded by it, and for such a
writer the bound of rule~6 is the remedy.

\begin{schema}{WOweOk}
  WrPc
\where
  pc~aWr = wowe \\
  current = fset \\
  ended = nocause \\
  hubreq = fclear \\
  ledger' = fset \\
  lk' = lk \\
  wire' = wire \\
  wcs' = wcs \\
  pc' = pc \oplus \{ aWr \mapsto wunl \}
\end{schema}

\begin{schema}{WOweRefused}
  WrPc
\where
  pc~aWr = wowe \\
  current = fset \\
  (ended \neq nocause \lor hubreq = fset) \\
  ledger' = ledger \\
  lk' = fclear \\
  wire' = wire \\
  wcs' = cdisp \\
  pc' = pc \oplus \{ aWr \mapsto wfail \}
\end{schema}

CONTROL: \texttt{owe\_ack} refuses on \texttt{closing}, a Hub request or a recorded
end, though the bytes are on the wire; the refusal is a \texttt{ConnectionClosedError},
which the writer's failure path reports as the display's.

\begin{schema}{WOweLate}
  WrLate
\where
  pc~aWr = wowe \\
  current = fclear \\
  a1late' = fset \\
  s2blk' = s2blk \\
  s2dec' = s2dec \\
  s5blk' = s5blk \\
  s2lat' = s2lat \\
  ledger' = ledger \\
  lk' = lk \\
  wire' = wire \\
  wcs' = wcs \\
  pc' = pc \oplus \{ aWr \mapsto wunl \}
\end{schema}

\begin{schema}{WUnlock}
  WrPc
\where
  pc~aWr = wunl \\
  lk' = fclear \\
  wire' = wire \\
  ledger' = ledger \\
  wcs' = wcs \\
  pc' = pc \oplus \{ aWr \mapsto pdone \}
\end{schema}

\texttt{owe\_ack}.  Once the bytes are on the wire the entry is made, whatever the
end state: a Hub request and a recorded cause do not refuse it (rule~7a).  Only a
dial already captured does, by \texttt{WOweLate}.  That is an \emph{under}-charge,
not an over-charge: the scene is on the wire, and is in neither the ledger the
capture took nor, for a \texttt{hub} capture, the carried set, so a death that
the scene causes cannot name it.  Any capture lands there, drained or not, since
a capture takes the state lock and not the socket lock: a Hub request, the
settler's decision on a socket it found unreadable, the bound, a decode-error
end.  It is the fifth exception to A1 (flag \texttt{a1late}), bounded to one
scene per dial.  The bound holds because the write then raises
\texttt{ConnectionClosedError}, an \texttt{OSError}; the replicator's recovery
re-marks the batch's scenes dirty for every failure but a shutdown flush
(\texttt{SendRecovery.recover} and \texttt{\_repaint}), so the scene goes out
again on the next dial and, if it is the poison scene, is charged there.  The
attribution is late by one dial; it is not lost.  The writer then lets go of the
lock.  The design admits no state in which the scene is on the wire of a
\emph{current} dial and not in the ledger (S6).

\begin{schema}{WStall}
  WrPc
\where
  pc~aWr = wsend \\
  lk' = lk \\
  wire' = wire \\
  ledger' = ledger \\
  wcs' = wcs \\
  pc' = pc \oplus \{ aWr \mapsto wstuck \}
\end{schema}

\begin{schema}{WStuckOut}
  WrPc
\where
  pc~aWr = wstuck \\
  current = fclear \\
  lk' = fclear \\
  wire' = wire \\
  ledger' = ledger \\
  wcs' = cdisp \\
  pc' = pc \oplus \{ aWr \mapsto wfail \}
\end{schema}

A writer that keeps making partial progress, so that its send timeout never
fires, and holds the lock for ever: a behaviour of the environment, like
\texttt{Block}.  It is released only by a capture, which shuts the socket and
makes \texttt{sendall} raise.  Such a writer is what the ceiling and the bound
exist for: without \texttt{SDrainLate} a settler with no refused asker behind it
waits for a writer that never lets go (the control \texttt{drain\_no\_late}).
The write then fails on a dial that is already captured, and its failure is
left alone (\texttt{WRecGone}).

\begin{schema}{WRec}
  WrRec
\where
  pc~aWr = wfail \\
  current = fset \\
  ((ended = nocause \land hubreq = fclear \land ended' = wcs) \lor
    (\lnot (ended = nocause \land hubreq = fclear) \land ended' = ended)) \\
  ((ended = nocause \land hubreq = fclear \land wcs = cdisp \land eof = fclear \land
      s5bad' = fset) \lor
    (\lnot (ended = nocause \land hubreq = fclear \land wcs = cdisp \land
      eof = fclear) \land s5bad' = s5bad)) \\
  slot' = slot \\
  debt' = debt \\
  current' = current \\
  sock' = sock \\
  capN' = capN \\
  hubreq' = hubreq \\
  s2bad' = s2bad \\
  s3bad' = s3bad \\
  s4bad' = s4bad \\
  pc' = pc \oplus \{ aWr \mapsto psu \}
\end{schema}

\begin{schema}{WRecGone}
  OnlyPc
\where
  pc~aWr = wfail \\
  current = fclear \\
  pc' = pc \oplus \{ aWr \mapsto pdone \}
\end{schema}

\texttt{\_write\_failed}.  The classification records the writer's cause with the
first-wins record, \emph{unless a Hub request is pending}: then the failure is the
Hub's, no cause is recorded (rule~7b), and the settler decides.  A display that
died is still found, by the settler's drain, which sees its end of file.  The
writer then asks for a settle, as every end does.  A cause recorded for a display
that is alive (\texttt{wcs = cdisp} with \texttt{eof = fclear}, which only a
failure that the display did not cause can produce) is the violation of S5; the
design cannot reach it, because \texttt{WSendDead} requires \texttt{eof = fset}
and \texttt{eof} never clears.

\section{The inline reader}

\begin{schema}{IAcq}
  DialPc
\where
  pc~aIn = pidle \\
  slot = sfree \\
  slot' = sinline \\
  ended' = ended \\
  debt' = debt \\
  current' = current \\
  sock' = sock \\
  capN' = capN \\
  hubreq' = hubreq \\
  pc' = pc \oplus \{ aIn \mapsto iread \}
\end{schema}

\begin{schema}{IReadIdle}
  OnlyPc
\where
  pc~aIn = iread \\
  blocked = fclear \\
  pc' = pc \oplus \{ aIn \mapsto ilast \}
\end{schema}

\begin{schema}{IReadNotice}
  OnlyPc
\where
  pc~aIn = iread \\
  blocked = fclear \\
  (ended \neq nocause \lor hubreq = fset \lor eof = fset) \\
  pc' = pc \oplus \{ aIn \mapsto pset \}
\end{schema}

\begin{schema}{ILastSet}
  OnlyPc
\where
  pc~aIn = ilast \\
  current = fset \\
  (ended \neq nocause \lor hubreq = fset) \\
  pc' = pc \oplus \{ aIn \mapsto pset \}
\end{schema}

\begin{schema}{ILastSkip}
  OnlyPc
\where
  pc~aIn = ilast \\
  (current = fclear \lor (ended = nocause \land hubreq = fclear)) \\
  pc' = pc \oplus \{ aIn \mapsto prel \}
\end{schema}

\begin{schema}{IBad}
  CapCarve
\where
  pc~aIn = iread \\
  blocked = fclear \\
  current = fset \\
  ((ended = nocause \land ended' = cdisp) \lor (ended \neq nocause \land ended' = ended)) \\
  current' = fclear \\
  capN' = capN + 1 \\
  slot' = slot \\
  debt' = debt \\
  sock' = sock \\
  hubreq' = hubreq \\
  s2bad' = s2bad \\
  ((capN \geq 1 \land s3bad' = fset) \lor (capN = 0 \land s3bad' = s3bad)) \\
  s4bad' = s4bad \\
  s5bad' = s5bad \\
  s2blk' = s2blk \\
  ((ended' = cdisp \land sock > 0 \land s2dec' = fset) \lor
    (\lnot (ended' = cdisp \land sock > 0) \land s2dec' = s2dec)) \\
  s5blk' = s5blk \\
  s2lat' = s2lat \\
  a1late' = a1late \\
  pc' = pc \oplus \{ aIn \mapsto prel \}
\end{schema}

A caller that holds the slot reads, notices an end (or a request) and owes a
settle, or finishes its read; its last check, \texttt{\_settle\_if\_ended}, then
either finds an end and owes a settle or finds none and releases.  Between the
last check and the release, anyone may end the dial: that is the window the
debt closes.  A frame that breaks the stream (\texttt{IBad}) records the
display's end and captures at once, undrained: the first carve-out.

\section{The listener thread}

\begin{schema}{RNotice}
  OnlyPc
\where
  pc~aRd = rtop \\
  blocked = fclear \\
  (hubreq = fset \lor ended \neq nocause \lor eof = fset) \\
  pc' = pc \oplus \{ aRd \mapsto pset \}
\end{schema}

\begin{schema}{RBad}
  CapCarve
\where
  pc~aRd = rtop \\
  blocked = fclear \\
  current = fset \\
  ((ended = nocause \land ended' = cdisp) \lor (ended \neq nocause \land ended' = ended)) \\
  current' = fclear \\
  capN' = capN + 1 \\
  slot' = slot \\
  debt' = debt \\
  sock' = sock \\
  hubreq' = hubreq \\
  s2bad' = s2bad \\
  ((capN \geq 1 \land s3bad' = fset) \lor (capN = 0 \land s3bad' = s3bad)) \\
  s4bad' = s4bad \\
  s5bad' = s5bad \\
  s2blk' = s2blk \\
  ((ended' = cdisp \land sock > 0 \land s2dec' = fset) \lor
    (\lnot (ended' = cdisp \land sock > 0) \land s2dec' = s2dec)) \\
  s5blk' = s5blk \\
  s2lat' = s2lat \\
  a1late' = a1late \\
  pc' = pc \oplus \{ aRd \mapsto rdone \}
\end{schema}

The listener's loop tests, in turn, the Hub's request (its stop flag), a
recorded end, and the display's end of file, and on any of them settles through
the same two steps that every other holder uses.  It records no cause of its
own and reads none twice.  A malformed frame is the first carve-out
(\texttt{RBad}).

\section{Refinement of crash\_attribution\_inflight\_ack\_drain.tex}

The abstract model that this specification refines is
\texttt{crash\_attribution\_inflight\_ack\_drain.tex}, the variant of
\texttt{crash\_attribution.tex} in which every \texttt{Discover*} step takes each
acknowledgement in the socket and in the reader's buffer into \texttt{acked} in
the same step (\texttt{Recv} and \texttt{Deliver} folded into the capture), and
in which a drain that applies an ack clears the carried set.  Its \texttt{ended}
takes \texttt{edisp}, \texttt{eunr} or \texttt{ehub}.  The gluing, concrete to
abstract:

\begin{itemize}
  \item \textbf{Cause.}  \texttt{cdisp}, \texttt{cunresp}, \texttt{chub} map to
    \texttt{edisp}, \texttt{eunr}, \texttt{ehub}.  A capture here is the abstract
    discovery of that cause.  The concrete \texttt{hubreq} is an abstract
    stutter: no abstract state changes when the Hub merely asks.  The abstract
    \texttt{HubEnd} is placed at the settle's observation (the drain), not at the
    ask.  That placement is the linearisation that makes the S5 residual a
    display that died \emph{after} the abstract \texttt{HubEnd}, that is before
    the abstract \texttt{CrashFrame}.
  \item \textbf{The drain.}  The concrete guard ``no ack unread when
    \texttt{cdisp} is captured'' (\texttt{sock = 0} at \texttt{SDecide}) refines
    the abstract in-step drain of the variant.  \texttt{sock} conflates the socket
    and the reader's undelivered buffer.  That is sound only because rule~1 makes
    the settler the owner of the frame buffer: a non-owner cannot empty what the
    reader has taken, but the owner's drain takes it to the ledger before it
    returns, so for the owner the two are one place.  This is the strongest
    argument for rule~1.  The abstract \texttt{Recv} then \texttt{Deliver} is one
    concrete step because the settler runs both without releasing the slot.
  \item \textbf{The gluing is an inclusion, not an equality of state.}  The
    unresponsive re-drain reads acknowledgements that the abstract model says
    were never sent: after \texttt{Unresponsive} a \texttt{Swap} adds nothing to
    the abstract \texttt{flight}.  The concrete \texttt{acked} is therefore a
    superset of the abstract one, and the relation that holds is
    \emph{concrete charges $\subseteq$ abstract charges}.  The inclusion gives A2.
    A1 holds for a different reason: the only acknowledgements the concrete
    settle reads beyond the abstract ones are, by T11, of scenes that survived a
    frame, so the poison scene is never among them.  The abstract model charges
    a scene that the concrete one does not only when a display recorded
    unresponsive later swaps and acknowledges that scene, and such a scene was
    not poison.
  \item \textbf{The obligations.}  A3 (a Hub-ended dial charges nothing) holds by
    the cause mapping, and by S3 (one capture).  A4 (a death is charged once)
    holds by S3.  A1 (an observed crash is charged) needs S1 and S5.  A2 (no
    scene outside the ledger is charged) needs S2.  O2's drain clause is S2.  All
    five hold on the holder's settle, with the re-drain of rule~4; without it A2
    and O2 fail (\texttt{split\_settle}).
\end{itemize}

\paragraph{Six exceptions to A1 and A2.}
The carve-outs are not bounds that the abstract model already contains.
\texttt{crash\_attribution.tex} has no charging stream-error end: its reply skew
is \texttt{HubEnd}, and every display or unresponsive discovery is guarded on an
empty flight and an empty reader buffer, so it never charges a scene whose ack
is unread.  Six concrete paths have no abstract counterpart.  On each, A1 or A2
fails by design; they are deliberate weakenings, not refinements.
\begin{enumerate}
  \item \texttt{s2dec}: a display charge with an ack unread, from a frame that
    broke the stream.  A2 fails.  Bound: the acks behind the bad frame, once.
  \item \texttt{s2blk}, display cause: a display charge with an ack unread,
    taken by the bound of rule~6.  A2 fails.  Bound: every ack queued while the
    holder stalled, including the holder's undelivered buffer, after the stop
    join.  The over-charge is accepted (DES-102).
  \item \texttt{s2blk}, unresponsive cause: the bound's capture of an
    unresponsive end, undrained.  A2 fails.  The bound is that of item~2.  The
    abstract model's B3 makes the retire instant the capture, but it still
    guards the charge on an empty flight and buffer.
  \item \texttt{s5blk}: \texttt{hub} recorded by the bound for a display that has
    already died.  An abstract \texttt{HubEnd} is disabled once the display is
    down.  The concrete result is that a poison death is carried and not charged,
    and a later ack read clears it.  A1 is lost by the B4 route, once, and only
    after the bound.
  \item \texttt{s2lat} (round 4): a display or unresponsive end captured, undrained,
    by the settler's own late capture, because a writer held the socket lock
    past the ceiling while an ack or the end of file was readable.  A2 fails.
    Bound: the acks readable at that moment, after the ceiling and the bound
    (the same ceiling-plus-bound as the settle that ends it); the over-charge is
    accepted, like item~2.
  \item \texttt{a1late} (round 4): a scene written to the wire while a capture
    lands between \texttt{sendall} returning and the owed-ack entry, so that it
    is in neither the captured ledger nor, for a \texttt{hub} capture, the
    carried set.  An \emph{under}-charge: A1 fails.  Bound: one scene per dial,
    and only for a capture in that window.  The write raises, the replicator's
    recovery re-marks the scene, and it goes out on the next dial, where a death
    it causes is charged; the attribution is late by one dial, not lost.
\end{enumerate}

This section is the refinement obligation list.  The matching note belongs in
\texttt{crash\_attribution.tex} next to its settle, and in the Bounds of DES-102;
the first is outside this mission's write set, so it is not edited here.

\section{Verification}

Commands (each run alone; MAX\_INITIALISATIONS must be at least 4: the slot
born with the listener or free, an unread acknowledgement or none).  Search mode
is stated with every count, because the breadth-first count and the default
count differ.

\begin{verbatim}
  fuzz docs/dial_settlement.tex                       (clean)
  P="-p MAX_OPERATIONS 100000 -p MAX_INITIALISATIONS 100"
  probcli docs/dial_settlement.tex -model_check -coverage $P          (deadlock, default mode)
  probcli docs/dial_settlement.tex -model_check -nodead -bf $P \
    -goal "s2bad = fset or s3bad = fset or s4bad = fset or s5bad = fset"
  S1a (timer-free; pwait counts as waiting, a stalled holder does not count):
  probcli docs/dial_settlement.tex -model_check -nodead -bf $P \
    -goal "current = fset & (ended /= nocause or hubreq = fset) &
           !a.(a : ACTOR => pc(a) : {pidle,pdone,rdone,rnone,pwait})"
  S1b (a stalled holder counts as stuck; a waiting asker is not):
  probcli docs/dial_settlement.tex -model_check -nodead -bf $P \
    -goal "current = fset & (ended /= nocause or hubreq = fset) &
           !a.(a : ACTOR => (pc(a) : {pidle,pdone,rdone,rnone} or
                (blocked = fset & pc(a) : {rtop,iread})))"
  The carve-outs are reachable, one flag each:
  probcli ... -goal "s2blk = fset"   / "s2dec = fset"   / "s5blk = fset"
  and the bound's undrained capture of an unresponsive end (carve-out 2):
  probcli ... -goal "s2blk = fset & ended = cunresp"
\end{verbatim}

\paragraph{Results for the design.}
Deadlock check, default search mode: 90906 states, 430865 transitions, no
counter example, \texttt{deadlocked:0}, 36 of 36 operations covered.  The three
goals, each \texttt{-bf} (breadth first, 90905 states analysed, 26 levels):
\texttt{s2bad or s3bad or s4bad or s5bad}, S1a and S1b are not found.  The
carve-out flags are found: \texttt{s2blk} in 4 steps (690 states; the shortest
trace is an unresponsive end, \texttt{DStartUnresp, RecEnd, SURefuse,
Timeout}), \texttt{s2dec} in 2 (66 states), \texttt{s5blk} in 5 (1106 states).
\texttt{s2blk = fset \& ended = cunresp} is found in the same 4 steps.  These
are expected, and they say that the carve-outs are real and bounded, not that
the design fails.

\paragraph{Errata (round 3, section ``The carve-outs of S2'').}
Round 2 flagged \texttt{s2blk} only when the bound captured a display end, and
said of the unresponsive end that a drain begun after the record precedes the
capture.  That holds for the holder's settle and is false for the bound's
capture, which is undrained for both causes.  \texttt{s2blk} now flags both, and
the state space grew from 86034 to 90906 states for that reason alone (the flag
is part of the state).  The checked properties are unchanged.

\paragraph{The goals, and which witness is faithful.}
S1a counts a waiting asker as quiescent and a stalled holder as in flight; S1b
counts a stalled holder as stuck and a waiting asker as in flight.  Both count an
actor that has not started (\texttt{pidle}) as quiescent, which makes a goal
witness \emph{over-report}: a later \texttt{close} could rescue the dial, as it
could in the code.  In the controls the witnesses of S1a and S1b are therefore
scenarios in which nothing running settles the dial, not proofs that nothing
ever will; the faithful witness of the wedge is the trace read against the code,
and where a deadlock is found (every actor counted, no timer: \texttt{defer\_only})
it is the one no modelled party can rescue.  With the timer present the
controls' deadlocks disappear, because the bound rescues them; that is why the
controls are checked by the S1a goal.

\paragraph{Controls, with search mode.}
Every goal run is \texttt{-nodead -bf} with the goal named; the deadlock run is
\texttt{-bf} with the default check.  ``Depth'' is the length of the trace
without the two set-up steps.
\begin{center}
\begin{tabular}{llll}
control & goal & states analysed & depth \\
start\_fail (docs) & S1a & 4627 & 7 \\
stop\_skip\_drain (docs) & \texttt{s2bad} & 19562 & 10 \\
inline\_strand (docs) & S1a & 3182 & 7 \\
split\_settle & \texttt{s2bad} & 9073 & 9 \\
atomic\_refuse & S1a & 5645 & 8 \\
atomic\_release & S1a & 5932 & 8 \\
defer\_only & deadlock & 6849 & 9 \\
timer\_hub\_only & S1b & 771 & 5 \\
s3\_unguarded & \texttt{s3bad} & 5552 & 8 \\
s4\_overwrite & \texttt{s4bad} & 1502 & 6 \\
hub\_record & \texttt{s5bad} & 123 & 3 \\
hub\_peek & \texttt{s5bad} & 166 & 3 \\
\end{tabular}
\end{center}
The ack that hides the end of file from a peek need not be in the initial
state: a living display acks and then dies (\texttt{Ack}, \texttt{Die}).  The
peek rule is therefore unsound whenever a display can have an ack in flight,
which is always; round 1's remark that an empty buffer made the peek sound was an
artefact of an ack that could not arrive, and is withdrawn here.

\paragraph{Round 4: the writer.}
The design now has the writer actor, the lock and the ledger (45 operations).
Same commands as above, with the S6 predicate added to the safety goal:
\begin{verbatim}
  P="-p MAX_OPERATIONS 200000 -p MAX_INITIALISATIONS 100"
  deadlock:  probcli docs/dial_settlement.tex -model_check -coverage $P
  safety:    ... -nodead -bf $P -goal "s2bad = fset or s3bad = fset or
               s4bad = fset or s5bad = fset or
               (wire = fset & ledger = fclear & lk = fclear & current = fset)"
  S1a, S1b:  the goals above, unchanged
\end{verbatim}
Results for the design (state space 1371734, up from 90906; peak 3.7 GB): the
deadlock check visits all states, 7384065 transitions, \texttt{deadlocked:0},
45 of 45 operations covered.  The safety goal (\texttt{s2bad}, \texttt{s3bad},
\texttt{s4bad}, \texttt{s5bad}, or S6 violated), S1a and S1b are each not found
over all 1371733 states, breadth first to level 30.  The writer controls:
\begin{center}
\begin{tabular}{llll}
control & goal & states analysed & depth \\
writer\_records\_display & \texttt{s5bad} & 3916 & 6 \\
writer\_records\_display & S6 predicate & 1629 & 5 \\
drain\_gives\_up & \texttt{s2bad} & 4645 & 6 \\
\end{tabular}
\end{center}
The traces.  \texttt{writer\_records\_display}, S5: \texttt{SPStart, WLock,
HubReq, WSendDone, WOweRefused, WRec}: a Hub request lands while a scene write is
in \texttt{sendall}; the write completes, the owed ack is refused, and the
failure records \texttt{display} for a display that is alive.  S6: a recorded end
(the keepalive's \texttt{unresponsive}) lands between \texttt{WSendDone} and the
owed ack, with the same result: the scene is on the wire of a current dial and
not in the ledger.  \texttt{drain\_gives\_up}, S2: \texttt{WLock, DStartUnresp,
RecEnd, SUGrant, SDrain, SDecide}: the keepalive's unresponsive end is recorded
while a writer holds the lock, the settle's drain finds the lock held and
reports nothing more with an ack readable, and the capture charges the scene
the display had acknowledged.  The design, with the guard \texttt{lk = fclear}
on \texttt{SDrain} and the owed ack made unconditional on a current dial, has
neither trace.  The counts are breadth-first, as everywhere in this section.

The two writer controls were checked in separate runs; the two small control
runs of the drain and of S6 happened to be issued together, and neither shares a
file with the other.

\paragraph{What is not modelled.}
A second discoverer, a second inline reader, more than one unread acknowledgement,
and the Hub's own stream-error end (\texttt{end\_stream(hub)}, a reply nobody was
owed): it records \texttt{hub} at once and captures undrained, because nothing can
be read past a stream already skewed; a \texttt{hub} capture charges nothing, so
the drain is moot.  The decode-error ends that record \texttt{display}
(\texttt{IBad}, \texttt{RBad}) are modelled, as the first carve-out.  The design
needs the refusal and its debt, the release and its debt-read, and the post-drain
record each to be one critical section; the three controls above show the
necessity, and the sufficiency is an implementation obligation.

Three further obligations lie beyond the model's actors, which are always
distinct threads and always same-host.
\begin{enumerate}
  \item \emph{The asker on the holder's thread does not wait.}  An interaction
    handler runs on the listener thread; a write that fails there reaches
    the slot's ask (\texttt{ask\_settle}, whose refusal hands back a
    \texttt{BoundedAsk}) with the slot held by its own thread.  Waiting out the
    bound on the thread that would settle captures undrained after two seconds,
    where returning would have drained and settled one step later.  Its refusal's
    debt is enough: it is the holder and settles on return.  The wait runs holding
    no lock that the holder's settle takes (the socket lock, \texttt{st}, the
    slot's condition).  The worst case of the waits on a Hub end is the bound,
    then the join: \texttt{ListenerThread.stop} takes up to both.
  \item \emph{The end of file the drain observes includes a reset and a TLS
    end.}  On a same-host \texttt{AF\_UNIX} socket a peer that closes with data
    \emph{it has not read} (a Hub frame still in the display's receive queue)
    gives the survivor \texttt{ECONNRESET} after the queued bytes, not an end of
    file.  A byte queued toward the Hub is ordinary data followed by a plain end
    of file.  A drain that cuts short on a reset, or on a TLS close, must say
    that an end was seen; otherwise S5 fails in the code with the model green.
  \item \emph{The complete drain is a same-host claim.}  An RST from a process
    that dies with unread input may discard what was queued, so on TCP a dead
    display's drain is not complete.  The claim that the drain is complete is
    restricted to same-host sockets, and that restriction is the contract; a
    cross-host end is outside it.
\end{enumerate}

\end{document}
